\documentclass[10pt,a4paper]{article}

\usepackage[utf8]{inputenc}
\usepackage[T1]{fontenc}
\usepackage{babel}
\usepackage[margin=1.6cm, top=1.4cm, bottom=1.4cm, headheight=14pt]{geometry}
\usepackage{xcolor}
\usepackage{enumitem}
\usepackage{titlesec}
\usepackage{fancyhdr}
\usepackage{graphicx}
\usepackage{tabularx}
\usepackage{booktabs}
\usepackage{microtype}
\usepackage{helvet}
\usepackage{listings}
\renewcommand{\familydefault}{\sfdefault}

\definecolor{accent}{RGB}{30, 80, 162}
\definecolor{muted}{RGB}{120, 120, 120}
\definecolor{rule}{RGB}{210, 215, 225}
\definecolor{dhlred}{RGB}{212, 0, 42}
\definecolor{codebg}{RGB}{245, 246, 250}

\lstset{
  basicstyle=\ttfamily\footnotesize,
  backgroundcolor=\color{codebg},
  frame=single,
  rulecolor=\color{rule},
  breaklines=true,
  columns=fullflexible,
  keepspaces=true,
  showstringspaces=false
}

\setlength{\parskip}{3pt}
\setlength{\parindent}{0pt}

\pagestyle{fancy}
\fancyhf{}
\fancyhead[L]{\small\color{muted} Système de Gestion de Flotte DHL \enspace\textbar\enspace ADR Semaine 1}
\fancyhead[R]{\small\color{muted} 14 mars 2026}
\fancyfoot[C]{\small\color{muted}\thepage}
\renewcommand{\headrulewidth}{0pt}

\titleformat{\section}
  {\normalsize\bfseries\color{accent}}{}{0em}{}[\ \textcolor{rule}{\rule[0.5ex]{0.3\linewidth}{0.4pt}}]
\titlespacing{\section}{0pt}{12pt}{6pt}

\titleformat{\subsection}
  {\small\bfseries\color{muted}}{}{0em}{}
\titlespacing{\subsection}{0pt}{6pt}{3pt}

\newcommand{\ruledline}{\textcolor{rule}{\rule{\linewidth}{0.4pt}}\vspace{4pt}}
\newcommand{\adrstatus}[1]{\hfill{\small\color{muted}\textit{Statut :} \textbf{#1}}}

\begin{document}

% ---- PAGE DE GARDE ----
\begin{titlepage}
  \centering
  \vspace*{2cm}
  {\LARGE\bfseries UNIVERSITÉ DE ROUEN NORMANDIE\\[4pt]
  FACULTÉ DES SCIENCES ET TECHNIQUES\par}
  \vspace{1.5cm}
  {\large ANNÉE UNIVERSITAIRE 2025--2026\par}
  \vspace{0.5cm}
  {\large Master 1 -- Génie Informatique et Logiciel\par}
  \vspace{1.5cm}
  \textcolor{rule}{\rule{\textwidth}{1.5pt}}
  \vspace{1cm}
  {\huge\bfseries\color{accent} Architecture Decision Records\par}
  \vspace{0.4cm}
  {\Large Système de Gestion de Flotte\par}
  \vspace{0.3cm}
  {\large\color{dhlred}\bfseries Contexte applicatif : DHL Express\par}
  \vspace{0.4cm}
  {\large Semaine 1\par}
  \vspace{1cm}
  \textcolor{rule}{\rule{\textwidth}{1.5pt}}
  \vfill
  {\small 14 mars 2026 \quad\textbar\quad Encadrants : Lydia \& Luc}
\end{titlepage}

\tableofcontents
\newpage

% ---- EN-TÊTE ----
{\fontsize{16}{20}\selectfont\bfseries\color{accent} Architecture Decision Records}\\[2pt]
{\large Système de Gestion de Flotte --- DHL Express}\\[4pt]
\ruledline

{\small Ce document recense les décisions d'architecture prises en semaine 1 pour le système de \textbf{suivi en temps réel de la flotte de véhicules DHL Express}. Le système couvre la gestion des livreurs, l'affectation des véhicules aux tournées, la maintenance du parc et les alertes automatiques. Les utilisateurs sont les \textbf{administrateurs}, \textbf{managers}, \textbf{techniciens} et \textbf{utilisateurs}.}
\vspace{4pt}

\begin{minipage}[t]{0.30\linewidth}
  \includegraphics[width=0.95\linewidth]{C1.png}
  \centering{\tiny\color{muted} Vue système (C1)}
\end{minipage}
\hfill
\begin{minipage}[t]{0.65\linewidth}
  \includegraphics[width=0.95\linewidth]{C2.png}
  \centering{\tiny\color{muted} Architecture des conteneurs (C2)}
\end{minipage}

\vspace{6pt}
\ruledline

% ================================================================
\section{ADR-001 --- Architecture microservices \adrstatus{Accepté}}

\subsection{Contexte}
Le système DHL doit gérer simultanément des responsabilités très différentes : le suivi GPS en temps réel des camionnettes de livraison, la gestion administrative des véhicules et des livreurs, la planification des maintenances, et l'envoi d'alertes aux managers. Ces besoins évoluent à des rythmes différents et ont des contraintes de performance distinctes. Une architecture monolithique rendrait difficile la scalabilité ciblée (par exemple, augmenter uniquement la capacité GPS en période de fêtes) et le déploiement indépendant de chaque partie.

\subsection{Décision}
Le système est découpé en \textbf{5 microservices indépendants}, chacun propriétaire de ses propres données et déployable séparément via Kubernetes :

\vspace{4pt}
\begin{tabularx}{\linewidth}{llX}
  \toprule
  \small\textbf{Service} & \small\textbf{Technologie} & \small\textbf{Responsabilité} \\
  \midrule
  Véhicules    & Spring Boot & CRUD du parc, gestion des statuts et de l'historique \\
  Conducteurs  & Spring Boot & Profils livreurs, permis de conduire, assignations \\
  Localisation & Node.js     & Streaming GPS temps réel, stockage TimescaleDB \\
  Maintenance  & Spring Boot & Planification et suivi des interventions techniques \\
  Événements   & Spring Boot & Analyse des règles métier, génération des alertes \\
  \bottomrule
\end{tabularx}

\vspace{4pt}
\subsection{Conséquences}
\textbf{Avantages :} isolation des pannes (une panne du service GPS n'impacte pas la gestion des conducteurs), scalabilité ciblée, déploiements indépendants par service, technologie adaptée à chaque besoin (polyglotte).

\noindent\textbf{Inconvénients :} complexité opérationnelle accrue au démarrage, nécessité d'une API Gateway et d'un bus d'événements pour la communication inter-services. Surcoût acceptable au regard des bénéfices à long terme.

% ================================================================
\section{ADR-002 --- gRPC pour le streaming GPS \adrstatus{Accepté}}

\subsection{Contexte}
Les boîtiers GPS embarqués dans les camionnettes DHL émettent leur position toutes les 5 secondes. Pour une flotte de 200 véhicules, cela représente environ 40 messages par seconde en continu. Le protocole REST/HTTP classique, basé sur un modèle requête-réponse, n'est pas adapté à ce flux unidirectionnel permanent et génèrerait une surcharge réseau importante.

\subsection{Décision}
Utiliser \textbf{gRPC (HTTP/2)} pour la communication entre les boîtiers GPS et le Service Localisation, ainsi qu'entre le Service Localisation et l'API Gateway pour l'affichage temps réel sur le dashboard manager. Les opérations classiques conservent REST ou GraphQL.

\vspace{4pt}
\begin{tabularx}{\linewidth}{lX}
  \toprule
  \small\textbf{Protocole} & \small\textbf{Usage} \\
  \midrule
  gRPC (server-streaming) & Positions GPS temps réel boîtier $\rightarrow$ Service Localisation $\rightarrow$ Frontend \\
  GraphQL                 & Requêtes agrégées du frontend (dashboard, historique) \\
  REST/JSON               & CRUD véhicules, conducteurs, maintenances \\
  \bottomrule
\end{tabularx}

\subsection{Conséquences}
\textbf{Avantages :} latence très faible grâce à la connexion persistante HTTP/2, sérialisation binaire (Protobuf) plus légère que JSON, adapté au streaming unidirectionnel continu.

\noindent\textbf{Inconvénients :} les navigateurs web ne supportent pas gRPC nativement — une couche de traduction \texttt{grpc-web} ou un proxy Envoy est nécessaire entre le frontend et le service. Complexité légèrement supérieure à REST pour les développeurs moins familiers.

% ================================================================
\section{ADR-003 --- Apache Kafka pour la communication asynchrone \adrstatus{Accepté}}

\subsection{Contexte}
Plusieurs actions dans le système DHL doivent déclencher des réactions dans d'autres services sans les bloquer : une alerte géofencing ne doit pas ralentir le traitement des positions GPS, et la fin d'une maintenance doit automatiquement remettre le véhicule en disponibilité. Une communication synchrone directe entre services créerait un couplage fort et des risques de cascade en cas de panne.

\subsection{Décision}
Utiliser \textbf{Apache Kafka} comme bus d'événements central. Chaque service publie des événements sur des topics dédiés et consomme ceux qui le concernent, sans rien savoir des autres services. Six topics sont définis pour le projet :

\vspace{4pt}
\begin{tabularx}{\linewidth}{lX}
  \toprule
  \small\textbf{Topic} & \small\textbf{Rôle} \\
  \midrule
  \texttt{flotte.vehicules.events}    & Changements d'état et de statut des véhicules \\
  \texttt{flotte.conducteurs.events}  & Cycle de vie des livreurs, expiration des permis \\
  \texttt{flotte.assignments.events}  & Affectation véhicule $\leftrightarrow$ livreur par tournée \\
  \texttt{flotte.localisation.gps}    & Positions GPS en temps réel (haute fréquence) \\
  \texttt{flotte.maintenance.events}  & Révisions planifiées, en cours ou terminées \\
  \texttt{flotte.alertes.events}      & Alertes consolidées vers le manager et le frontend \\
  \bottomrule
\end{tabularx}

\subsection{Conséquences}
\textbf{Avantages :} découplage total entre services, aucune perte de message (rétention configurable), rejeu possible en cas d'erreur, scalabilité horizontale via les partitions.

\noindent\textbf{Inconvénients :} un composant d'infrastructure supplémentaire à déployer et surveiller. La cohérence des données devient \textit{éventuelle} (eventual consistency) plutôt qu'immédiate, ce qui est acceptable pour les cas d'usage du projet.

% ================================================================
\section{ADR-004 --- Bases de données spécialisées par service \adrstatus{Accepté}}

\subsection{Contexte}
Le système manipule trois catégories de données aux caractéristiques très différentes. Les positions GPS s'accumulent à très haute fréquence et nécessitent des requêtes temporelles efficaces. Les données métier (véhicules, livreurs, maintenances) sont relationnelles et transactionnelles. Certaines données très consultées (état courant d'un véhicule, session utilisateur) doivent être accessibles en quelques millisecondes.

\subsection{Décision}
Appliquer le pattern \textbf{Database per Service} avec trois outils complémentaires :

\vspace{4pt}
\begin{tabularx}{\linewidth}{llX}
  \toprule
  \small\textbf{Outil} & \small\textbf{Type} & \small\textbf{Rôle dans le projet DHL} \\
  \midrule
  PostgreSQL   & Relationnel       & Véhicules, livreurs, maintenances, alertes (données métier) \\
  TimescaleDB  & Séries temporelles & Historique des positions GPS (\texttt{location\_readings}) \\
  Redis        & Cache en mémoire  & État courant des véhicules, sessions Keycloak, données fréquentes \\
  \bottomrule
\end{tabularx}

\vspace{4pt}
Le schéma PostgreSQL est composé de 6 tables principales : \texttt{vehicles}, \texttt{vehicle\_assignments}, \texttt{drivers}, \texttt{driver\_licenses}, \texttt{maintenance\_records} et \texttt{alerts}. Ce périmètre est volontairement limité pour garantir un MVP fonctionnel.

\vspace{4pt}
\begin{minipage}[t]{0.48\linewidth}
  \centering
  \includegraphics[width=0.80\linewidth]{b1.png}\\[6pt]
  \includegraphics[width=0.80\linewidth]{b3.png}
\end{minipage}
\hfill
\begin{minipage}[t]{0.48\linewidth}
  \centering
  \includegraphics[width=0.80\linewidth]{b2.png}\\[6pt]
  \includegraphics[width=0.80\linewidth]{b4.png}
\end{minipage}

\subsection{Conséquences}
\textbf{Avantages :} chaque outil est optimisé pour son type de données, pas de contention entre charges différentes, TimescaleDB réduit considérablement la taille des données GPS grâce à la compression automatique.

\noindent\textbf{Inconvénients :} trois systèmes à opérer et sauvegarder. Les jointures entre données de services différents ne sont pas possibles directement — elles sont réalisées au niveau applicatif (API Gateway / GraphQL).

% ================================================================
\section{ADR-005 --- Keycloak pour l'authentification et les droits \adrstatus{Accepté}}

\subsection{Contexte}
Les utilisateurs du système DHL n'ont pas tous accès aux mêmes fonctionnalités. Un manager peut affecter des tournées mais ne peut pas modifier les données de maintenance. Un technicien accède aux fiches d'intervention mais pas aux données des autres utilisateurs. Implémenter cette logique de droits dans chaque microservice serait coûteux à maintenir et source d'incohérences.

\subsection{Décision}
Centraliser l'authentification et la gestion des droits dans \textbf{Keycloak} (SSO open-source). Chaque requête passe par l'API Gateway qui valide le \textbf{JWT} émis par Keycloak avant de router vers le bon service.

\vspace{4pt}
\begin{tabularx}{\linewidth}{lX}
  \toprule
  \small\textbf{Rôle Keycloak} & \small\textbf{Droits dans le système DHL} \\
  \midrule
  \texttt{admin}       & Accès complet à tous les services et données, gestion des comptes \\
  \texttt{manager}     & Consultation globale, validation des maintenances, gestion des tournées et du parc \\
  \texttt{technicien}  & Lecture/écriture sur les fiches de maintenance, consultation des véhicules \\
  \texttt{utilisateur} & Consultation de ses propres données (véhicule assigné, tournée en cours) \\
  \bottomrule
\end{tabularx}

\vspace{4pt}
Les microservices ne gèrent \textbf{aucun mot de passe} directement. Ils se contentent de vérifier la signature du JWT et d'extraire le rôle pour appliquer les autorisations (RBAC).

\subsection{Conséquences}
\textbf{Avantages :} gestion des droits centralisée et cohérente, SSO (connexion unique pour toutes les parties du système), standard OAuth2/OIDC, interface d'administration Keycloak disponible pour les ops.

\noindent\textbf{Inconvénients :} Keycloak est un point critique --- s'il est indisponible, aucune authentification n'est possible. Il doit donc être déployé en mode haute disponibilité (minimum 2 réplicas Kubernetes).

% ================================================================
\section{ADR-006 --- Observabilité avec OpenTelemetry \adrstatus{Accepté}}

\subsection{Contexte}
Dans une architecture microservices, une requête utilisateur traverse plusieurs services avant de retourner une réponse. En cas d'erreur ou de lenteur, il est très difficile d'identifier quel service est responsable sans outillage dédié. De plus, les métriques opérationnelles (CPU, mémoire, latence) et métier (nombre de véhicules actifs, alertes générées) sont indispensables au pilotage du système DHL.

\subsection{Décision}
Adopter \textbf{OpenTelemetry} comme standard unique d'instrumentation, couvrant les trois piliers de l'observabilité :

\vspace{4pt}
\begin{tabularx}{\linewidth}{llX}
  \toprule
  \small\textbf{Pilier} & \small\textbf{Outil} & \small\textbf{Usage dans le projet} \\
  \midrule
  Traces   & Jaeger     & Suivi distribué d'une requête à travers tous les services \\
  Métriques & Prometheus + Grafana & Tableaux de bord opérationnels et métier (véhicules actifs, alertes) \\
  Logs     & Loki       & Agrégation centralisée des logs structurés (JSON) de tous les services \\
  \bottomrule
\end{tabularx}

\vspace{4pt}
Chaque service est instrumenté via le SDK OpenTelemetry de son langage et exporte vers un \textbf{OTel Collector} central, qui redistribue ensuite vers Jaeger, Prometheus et Loki. Un \texttt{correlation\_id} est propagé dans tous les messages Kafka et requêtes HTTP pour relier les traces bout-en-bout.

\subsection{Conséquences}
\textbf{Avantages :} standard ouvert (pas de vendor lock-in), instrumentation mutualisée entre tous les services, corrélation logs/traces/métriques facilitée, Grafana unifie la visualisation.

\noindent\textbf{Inconvénients :} légère surcharge mémoire et CPU liée à l'instrumentation. Configuration initiale du Collector et des exporters à prévoir en semaine 2.

% ================================================================
\vspace{8pt}
\ruledline
{\small\color{muted}
\textbf{Synthèse des décisions} --- Les 6 ADR ci-dessus définissent les fondations techniques du projet. L'approche retenue privilégie un \textbf{MVP fonctionnel} : périmètre de base de données volontairement simple (6 tables), stack éprouvée, et enrichissements progressifs semaine après semaine.
}

\end{document}
